\section{Deep Learning for Code Translation}

%SzB: Ez inkább theoretical foundations jelleg. Ami nem rossz egyáltalán, sőt
% Lehet át lehetne nevezni a címet theoretical foundations and related work-re
% De a related work rész innen hiányzik. Ha már kódfordítás, akkor itt fel kéne tüntetni, hogy pl. akkor RNN-eket ki használt fordításra. Meg s többbit mindent

The emergence of deep learning has significantly transformed the field of code
translation by enabling data-driven approaches that learn patterns directly from
large-scale datasets. These methods are largely inspired by advances in neural
machine translation (NMT), where sequence-to-sequence (seq2seq) models have
become a dominant paradigm.

Early deep learning approaches to code translation relied on encoder–decoder
architectures based on recurrent neural networks (RNNs), including Long
Short-Term Memory (LSTM) and Gated Recurrent Unit (GRU) models. In these
architectures, the encoder processes the input sequence and encodes it into a
fixed-length vector representation, which is then used by the decoder to
generate the target sequence \cite{seq2seq_original}. While effective in
modeling sequential data, these models suffer from limitations when handling
long sequences, as important information may be lost during encoding
\cite{seq2seq_limitations}. 

To address these limitations, the attention mechanism was introduced, allowing
the model to dynamically focus on relevant parts of the input sequence during
decoding. This significantly improved translation quality by overcoming the
bottleneck of fixed-length representations and enabling better handling of
long-range dependencies \cite{attention_bahdanau}. 

These neural machine translation techniques were later adapted to source code,
where programs are treated as sequences of tokens similar to natural language.
This adaptation enabled models to learn mappings between programming languages
without relying on handcrafted rules. However, unlike natural language, source
code has strict syntax and semantics, which introduces additional challenges for
deep learning models.

Despite the improvements brought by attention mechanisms, RNN-based seq2seq
models remain difficult to parallelize and scale efficiently. This limitation
led to the development of the Transformer architecture, which relies entirely on
attention mechanisms and eliminates recurrence, enabling more efficient training
and improved performance in sequence modeling tasks \cite{vaswani2017}. 

The introduction of Transformers marked a major shift in the field and laid the
foundation for modern approaches to code translation, which are discussed in the
following section.